\section{Architecture \& Implementation}
\label{method}

% Goal: describe what was built and why each design decision was made.
% This is an architecture/methods section, not a biology section.

% ── 1. Datasets ────────────────────────────────────────────────────────
%   TABLE — 5 datasets:
%   Dataset       | RBP    | Cell  | ENCODE accessions (IP rep1, IP rep2, SMInput)
%   RBFOX2_K562   | RBFOX2 | K562  | ENCFF537RYR, ENCFF296GDR, ENCFF212IIR
%   RBFOX2_HepG2  | RBFOX2 | HepG2 | ENCFF239CML, ENCFF170YQV, ENCFF515BTB
%   QKI_K562      | QKI    | K562  | ENCFF698BKX, ENCFF012WMS, ENCFF070RME
%   QKI_HepG2     | QKI    | HepG2 | ENCFF567ADV, ENCFF862YVK, ENCFF015GLL
%   PUM1_K562     | PUM1   | K562  | ENCFF064COB, ENCFF583QFB, ENCFF222HEX
%   All BAMs: GRCh38, deduplicated (ENCODE pipeline). chr21 fast mode
%   used throughout (~minutes/trial vs. hours genome-wide).

% ── 2. Evaluation pipeline ─────────────────────────────────────────────
%   Single shared evaluate_config() entry point called by both optimizers:
%   (1) Merge IP replicates (samtools).
%   (2) Run pureclip2 on merged BAM → binding_regions.bed.
%   (3) Run pureclip2 per-replicate → per-replicate regions.
%   (4) Post-process: filter by min_crosslink_events, force footprint
%       width → binding_sites.bed.
%   (5) Score → score_report.json.
%   Figure: pipeline flowchart.

% ── 3. Tunable parameters ──────────────────────────────────────────────
%   TABLE — 7 parameters with bounds and what they control:
%   PureCLIP:        bandwidth_nt [20–100], merge_distance_nt [4–16],
%                    high_precision_mode {0,1}, use_input_covariate {0,1}
%   Post-processing: min_crosslink_events [2–6], force_width [3–15],
%                    cluster_gap_width [4–16]
%   Justify why each is tunable vs. fixed (brief sentence per group).

% ── 4. Composite objective function ───────────────────────────────────
%   composite = (0.5·repro + 0.25·motif + 0.25·recall)
%               × min(1, n_sites / 10)
%   Justify weights: reproducibility = strongest evidence of genuine
%   binding; motif + recall = biological validity + known-site recovery.
%   Collapse guard: prevents degenerate solutions (1–2 "perfect" sites);
%   motivated by an observed Optuna failure (QKI collapsed to 1 site,
%   composite 0.78) before the guard was introduced.

% ── 5. LLM optimizer ──────────────────────────────────────────────────
%   Model: DeepSeek (deepseek-chat, T=0.2). Each iteration receives:
%   biological priors (protein, motif, footprint ~9 nt), search bounds,
%   current parameters, and a progress table with per-iteration score
%   deltas. Proposes one parameter change with reasoning.
%   Validation: JSON parsed, checked against bounds + prior trials;
%   up to 3 retries with rejection reason in prompt.
%   Stopping: max_iterations (10–12) or 10 steps without improvement.

% ── 6. Optuna optimizer ───────────────────────────────────────────────
%   TPESampler (seed 42), maximising composite_objective.
%   Budget: 12–15 trials. No domain knowledge; purely data-driven.
%   Same evaluate_config() call, same decisions.jsonl output format.
%   KEY DESIGN DECISION: identical evaluation for both optimizers ensures
%   the comparison isolates proposal strategy, not pipeline differences.
